% !TEX root = ../main.tex
\section{Implementation}\label{sec:implementation}

We describe the open-source implementation of \textsc{zk-eidas}, available under the Apache~2.0
licence at \url{https://zk-eidas.com}.
The system is written in Rust with a C++ proving core, a React~19 web frontend, and an Axum
HTTP API backend.
We structure the presentation as four subsections covering the crate architecture
(\S\ref{sec:impl-arch}), the circuit compilation pipeline (\S\ref{sec:impl-circuits}),
the TSP attestation layer (\S\ref{sec:impl-tsp}), and the deployment model
(\S\ref{sec:impl-deploy}).

%% ───────────────────────────────────────────────────────────────────
\subsection{Architecture Overview}\label{sec:impl-arch}

\paragraph{Rust workspace.}
\textsc{zk-eidas} is organised as a Cargo workspace of eight crates with clearly separated
responsibilities.
Table~\ref{tab:crates} summarises each crate.

\begin{table}[t]
  \centering
  \caption{Rust workspace crates and their responsibilities.}\label{tab:crates}
  \begin{tabular}{ll}
    \toprule
    Crate & Responsibility \\
    \midrule
    \texttt{zk-eidas-types}  & Witness, proof, and credential data structures \\
    \texttt{zk-eidas-mdoc}   & ISO/IEC~18013-5~\cite{iso18013} CBOR parsing, COSE signature extraction \\
    \texttt{zk-eidas-utils}  & Date conversion, CBOR encoding utilities \\
    \texttt{zk-eidas}        & Facade: escrow, predicate builder, proof orchestration \\
    \texttt{longfellow-sys}  & FFI bindings to Longfellow~\cite{longfellow} C++, circuit management \\
    \texttt{zk-eidas-wasm}   & WASM verifier (Ligero verification in pure Rust) \\
    \texttt{cbor-print}      & CBOR debugging and pretty-printing utility \\
    \texttt{demo-api}        & Axum HTTP API server (issuer, holder, verifier endpoints) \\
    \bottomrule
  \end{tabular}
\end{table}

\paragraph{FFI architecture.}
The interface between Rust and the Longfellow C++ library is encapsulated entirely within the
\texttt{longfellow-sys} crate.
At build time, \texttt{build.rs} invokes CMake to compile Longfellow's \texttt{mdoc\_static}
target, which links the sumcheck and Ligero backends into a static archive.
\texttt{bindgen} then generates typed Rust FFI declarations from a thin \texttt{wrapper.h} header.
Three functions constitute the public FFI surface visible to the rest of the workspace:
\begin{itemize}
  \item \texttt{generate\_circuit(spec, \&mut circuit, \&mut len)} --- compiles a \texttt{kZkSpecs}
        entry into binary circuit bytes and returns them via a heap-allocated buffer.
  \item \texttt{run\_mdoc\_prover(spec, circuit, attrs, mdoc, \ldots, \&mut proof)} --- executes
        the Longfellow prover and writes the proof together with the public outputs
        (nullifier hash, binding hash, escrow digest) into caller-provided buffers.
  \item \texttt{run\_mdoc\_verifier(spec, circuit, attrs, \ldots, proof)} --- verifies a proof,
        returning a typed error code on failure.
\end{itemize}
All three calls are wrapped by the \texttt{longfellow\_sys::safe} module, which converts raw
pointers into idiomatic Rust slices, performs bounds checks, and translates C error codes into
a typed \texttt{MdocError} enum.

\paragraph{End-to-end pipeline.}
A presentation flows through five stages:
\begin{enumerate}
  \item \emph{mdoc parsing}: \texttt{zk-eidas-mdoc} decodes the raw CBOR credential into
        structured form and extracts the COSE\_Sign1 envelope and MSO.
  \item \emph{Witness construction}: \texttt{zk-eidas} builds the attribute list (namespace,
        identifier, CBOR value, predicate type), escrow fields, contract hash, and issuer public
        key coordinates into a \texttt{ProveRequest}.
  \item \emph{Circuit lookup}: the API layer retrieves the pre-generated binary circuit matching
        the requested attribute count from the on-disk cache (see \S\ref{sec:impl-circuits}).
  \item \emph{Proving}: \texttt{run\_mdoc\_prover} executes the Sumcheck+Ligero prover and
        returns proof bytes together with the three public hashes.
  \item \emph{Verification}: \texttt{run\_mdoc\_verifier} re-runs the hash and signature
        sub-circuits against the supplied public inputs and rejects any mismatch.
\end{enumerate}

%% ───────────────────────────────────────────────────────────────────
\subsection{Circuit Specifications}\label{sec:impl-circuits}

\paragraph{Variant table.}
Longfellow circuits are parameterised by the number of attributes.
We compile four variants, indexed by \texttt{kZkSpecs[0..3]} (specification version~7):

\begin{table}[t]
  \centering
  \caption{Circuit variants by attribute count.
           Sizes are the on-disk binary blobs as used at runtime (uncompressed, written
           by \texttt{generate-circuits} and loaded by the API server).}\label{tab:circuits}
  \begin{tabular}{ccc}
    \toprule
    Attributes & \texttt{kZkSpecs} index & Cached file size \\
    \midrule
    1 & 0 & 337 KB \\
    2 & 1 & 349 KB \\
    3 & 2 & 365 KB \\
    4 & 3 & 384 KB \\
    \bottomrule
  \end{tabular}
\end{table}

\paragraph{Compilation cost and caching strategy.}
Generating a circuit from a \texttt{kZkSpecs} entry is computationally expensive: on our
benchmark hardware (Intel~2.60~GHz, 8~cores), circuit generation takes
${\approx}17$--$19$~seconds per variant and consumes ${\approx}1.3$--$1.4$~GB of RAM at peak.
We therefore pre-generate all four circuits \emph{once} during the Docker build (Stage~1a,
the \texttt{circuit-builder} image) and write the resulting binary blobs to an on-disk cache
directory.
At API server startup the cache is memory-mapped with a \texttt{tokio::sync::OnceCell},
reducing cold-start circuit load time to ${\approx}0.02$~ms per variant across all
four sizes (Table~\ref{tab:circuits}).
This separation ensures that a change to the application source triggers a rebuild of
Stage~1b (\texttt{rust-builder}) without invalidating the circuit cache layer.

\paragraph{DSL extensions.}
The Longfellow circuit DSL is extended at the boundary between the \texttt{longfellow-sys}
crate and the C++ backend with four computed outputs that are not part of the vanilla
Longfellow mdoc spec:
\begin{itemize}
  \item \emph{Predicate evaluation}: the four primitive predicate types
        (equality, $\leq$, $\geq$, $\neq$; see \S\ref{sec:predicates})
        are encoded as a \texttt{verify\_type} field on each \texttt{RequestedAttribute}
        and evaluated by the hash sub-circuit using a constant-time MUX gate.
  \item \emph{Nullifier}: $\eta = \mathsf{SHA\text{-}256}(m \| c)$ computed from the raw mdoc
        bytes and a caller-supplied 64-byte contract hash (\S\ref{sec:nullifiers}).
  \item \emph{Holder binding}: $\beta = \mathsf{SHA\text{-}256}(\mathsf{attr}_0[0..31])$
        derived from the first attribute value (\S\ref{sec:nullifiers}).
  \item \emph{Escrow digest}: $\delta = \mathsf{SHA\text{-}256}(e_0 \| \cdots \| e_7)$ over the
        eight zero-padded escrow field values (\S\ref{sec:escrow-construction}).
\end{itemize}
All four computations are wired into the \emph{same} hash sub-circuit evaluation pass,
adding only a constant number of SHA-256 compression calls to the circuit without increasing
circuit depth.

%% ───────────────────────────────────────────────────────────────────
\subsection{TSP Attestation Layer}\label{sec:impl-tsp}

To satisfy QEAA-grade (Qualified Electronic Attestation of Attributes) requirements under
eIDAS~2.0, the API server acts as a lightweight \emph{Trusted Service Provider} (TSP).
After a proof is verified by the Longfellow verifier, the holder can request a
\texttt{DataIntegrityProof} attestation from the TSP endpoint.

\paragraph{Attestation construction.}
The TSP builds a W3C Verifiable Credential document containing the proof CID, the verified
predicate description, the credential type, and a \texttt{verificationResult} field set to
\texttt{"valid"}.
The document body is serialised to canonical JSON and a SHA-256 digest is computed over it.
The TSP then signs the digest with ECDSA over NIST~P-256, producing a DER-encoded signature
that is hex-encoded into the \texttt{proofValue} field:
\begin{equation*}
  \sigma_{\mathrm{TSP}} \;=\; \mathsf{ECDSA}_{\mathit{sk}_{\mathrm{TSP}}}\!\bigl(\mathsf{SHA\text{-}256}(J)\bigr),
\end{equation*}
where $J$ is the canonical JSON serialisation of the VC body and $\mathit{sk}_{\mathrm{TSP}}$
is an ephemeral P-256 signing key generated at server startup.

\paragraph{Proof object.}
The \texttt{proof} object attached to the VC follows the \texttt{ecdsa-jcs-2019} cryptosuite
convention~\cite{w3c-di}:
\begin{itemize}
  \item \texttt{type}: \texttt{"DataIntegrityProof"}
  \item \texttt{cryptosuite}: \texttt{"ecdsa-jcs-2019"}
  \item \texttt{verificationMethod}: the uncompressed P-256 public key (hex-encoded)
  \item \texttt{proofValue}: the ECDSA signature bytes (hex-encoded)
\end{itemize}

\paragraph{Verification flow.}
A relying party that receives a TSP-attested credential bundle performs two independent checks:
(i)~it runs \texttt{run\_mdoc\_verifier} to confirm the Sumcheck+Ligero proof is valid against
the stated predicates and public hashes; and
(ii)~it verifies $\sigma_{\mathrm{TSP}}$ against the TSP's public key and the canonical JSON
of the VC body.
Only a bundle passing both checks is accepted as a QEAA-grade presentation.
The two-layer design allows the ZK verifier to remain stateless and crypto-agnostic, while
the TSP signature provides the regulatory non-repudiation guarantee.

%% ───────────────────────────────────────────────────────────────────
\subsection{Deployment}\label{sec:impl-deploy}

\paragraph{Docker multi-stage build.}
\textsc{zk-eidas} ships as a single Docker image built from a four-stage \texttt{Dockerfile}:

\begin{description}
  \item[\texttt{circuit-builder} (Stage~1a).] Compiles the \texttt{generate-circuits} binary and
    runs it to produce all four circuit cache files.
    This stage depends only on \texttt{vendor/longfellow-zk/} and the \texttt{crates/} directory;
    changes to application code do not invalidate this layer, so the ${\approx}70$-second
    circuit-generation step is paid only when the Longfellow vendor tree or the circuit spec
    changes.

  \item[\texttt{rust-builder} (Stage~1b, inherits from \texttt{circuit-builder}).] Compiles the
    \texttt{zk-eidas-demo-api} and \texttt{pre-warm} binaries against the already-linked
    Longfellow static library.

  \item[\texttt{wasm-builder} (Stage~1c, parallel).] Uses \texttt{wasm-pack} to build the
    \texttt{zk-eidas-wasm} crate into a WebAssembly module providing pure-Rust Ligero
    verification in the browser.

  \item[\texttt{web-builder} (Stage~2).] Runs \texttt{npm run build} for the React~19
    frontend (Vite bundler, TanStack Router, Tailwind~CSS) and consumes the WASM artefact
    from Stage~1c.

  \item[\texttt{runtime} (Stage~3).] A minimal Ubuntu~24.04 image containing only the API
    binary, the pre-generated circuit cache, the pre-warmed proof cache, and the compiled
    frontend static assets.
    The circuit files are copied from Stage~1a, so the runtime image carries no build
    toolchain and no Rust source.
\end{description}

\paragraph{Production deployment.}
The production instance runs at \url{https://zk-eidas.com} on Fly.io, behind an Nginx
reverse proxy managed by \texttt{supervisord}.
The API server and the static frontend are co-located in the same container.
A pre-warm binary is executed once after each deployment to populate the proof cache with
reference credentials, so that the demo endpoints respond without incurring a first-time
proving cost.

\paragraph{Licence and availability.}
All source code is released under the Apache~2.0 licence.
The Longfellow C++ library, vendored under \texttt{vendor/longfellow-zk/}, is governed by its
own Apache~2.0 licence~\cite{longfellow}.
